\section*{Bab \ref{ch:b2:nvs} --- Ruang Vektor Bernorma}

\begin{solution}{exo:b2:nvs:1}
Bola satuannya: belah ketupat ($\norm\cdot_1$), cakram ($\norm\cdot_2$),
dan persegi ($\norm\cdot_\infty$), yang bersarang menurut urutan itu.
Ketaksamaannya: $\norm x_\infty \leq \norm x_2$ (sebab satu kuadrat
paling banyak sama dengan jumlahnya); $\norm x_2 \leq \norm x_1$ (setelah
dikuadratkan: $x_1^2 + x_2^2 \leq (\abs{x_1} + \abs{x_2})^2$); dan
$\norm x_1 \leq 2\norm x_\infty$ (dua suku, masing-masing $\leq \max$).
Ketajamannya: $(1, 0)$ membuat dua yang pertama menjadi kesamaan;
sedangkan $(1, 1)$ membuat $\norm x_1 = 2\norm x_\infty$ dan sekaligus
menunjukkan bahwa $\norm x_2 = \sqrt2 \norm x_\infty$ dan $\norm x_1 =
\sqrt2 \norm x_2$ merupakan rasio ekstrem pada arah sebaliknya.
\end{solution}

\begin{solution}{exo:b2:nvs:2}
Aksioma \omterm{def:b2:nvs:norm}{normanya}: kehomogenan dan ketaksamaan segitiganya diwarisi suku
demi suku; pemisahannya: $N(f) = 0$ memaksa $f' = 0$ (jadi $f$ konstan)
sekaligus $f(0) = 0$, sehingga $f = 0$. Jadi ia memang \omterm{def:b2:nvs:norm}{norma}.

Perbandingannya: $\norm f_\infty \leq N(f)$, sebab $\abs{f(x)} \leq
\abs{f(0)} + \abs{\int_0^x f'} \leq \abs{f(0)} +
\norm{f'}_\infty$. Konversnya gagal: ambil $f_n(x) = \frac1n
\sin(nx)$: maka $\norm{f_n}_\infty \leq \frac1n \to 0$ padahal
$N(f_n) = 0 + \norm{\cos(nx)}_\infty = 1$. Jadi tak ada konstanta $C$
yang memberi $N \leq C\norm\cdot_\infty$.
\end{solution}

\begin{solution}{exo:b2:nvs:3}
$\abs{u(f)} \leq \norm f_\infty \int_0^1 \eu^t\,\dd t = (\eu -
1)\norm f_\infty$, dengan kesamaan untuk $f \equiv 1$: jadi $\vertiii u =
\eu - 1$.

Geseran: $\norm{S(x)}_\infty = \max(\abs{x_2}, \dots, \abs{x_n}) \leq
\norm x_\infty$, dengan kesamaan di $x = e_2$: jadi $\vertiii S = 1$
(untuk $n \geq 2$).
\end{solution}

\begin{solution}{exo:b2:nvs:4}
Batas atasnya: untuk $\norm x_\infty \leq 1$,
\[
\abs{(Ax)_i} = \Bigl|\sum_j a_{ij}x_j\Bigr| \leq \sum_j
\abs{a_{ij}},
\]
jadi $\norm{Ax}_\infty \leq \max_i \sum_j \abs{a_{ij}}$. Tercapainya:
misalkan $i_0$ mewujudkan maksimumnya lalu ambil $x_j =
\operatorname{sign}(a_{i_0 j})$ (unsur bermodulus $1$): maka $(Ax)_{i_0}
= \sum_j \abs{a_{i_0 j}}$. Karenanya rumusnya berlaku. Untuk matriks yang
diberikan: jumlah barisnya $3$ dan $4$, jadi $\vertiii A_\infty = 4$.
\end{solution}

\begin{solution}{exo:b2:nvs:5}
\emph{Deret geometri:} untuk $\vertiii B < 1$, deret $\sum
B^k$ konvergen mutlak di $\mathcal{M}_n(K)$ yang Banach
(\cref{thm:b2:nvs:absoluteconvergence}, sebab $\vertiii{B^k} \leq
\vertiii B^k$), dan
\[
(I - B)\sum_{k=0}^{K} B^k = I - B^{K+1} \longrightarrow I :
\]
jadi berkat kekontinuan hasil kalinya
(\cref{prop:b2:nvs:bilinear}), $(I - B)\sum_{k\geq0} B^k = I$, sehingga
$I - B$ punya invers berupa jumlah itu (dan
$\vertiii{(I-B)^{-1}} \leq \frac{1}{1 - \vertiii B}$).

\emph{Keterbukaannya:} untuk $A$ yang berinvers dan $\vertiii H <
\frac{1}{\vertiii{A^{-1}}}$: $A + H = A(I + A^{-1}H)$ dengan
$\vertiii{A^{-1}H} \leq \vertiii{A^{-1}}\vertiii H < 1$: jadi berinvers.
Karenanya sebuah bola di sekitar $A$ tetap berada di $GL_n$.

\emph{Kekontinuan pembalikannya:} dengan $B = -A^{-1}H$,
\[
(A + H)^{-1} - A^{-1} = \bigl((I - B)^{-1} - I\bigr)A^{-1}
= \Bigl(\sum_{k \geq 1} B^k\Bigr) A^{-1},
\]
yang \omterm{def:b2:nvs:norm}{normanya} $\leq \frac{\vertiii B}{1 - \vertiii B}\vertiii{A^{-1}}
\to 0$ ketika $H \to 0$.
\end{solution}

\begin{solution}{exo:b2:nvs:6}
\omterm{def:b2:metric:continuity}{Kontinu} $\Rightarrow$ kernelnya tertutup: sebab ia prapeta $\{0\}$ yang
tertutup (\cref{thm:b2:metric:globalcontinuity}).

Sebaliknya, andaikan $\ker\varphi$ tertutup dan $\varphi \neq 0$. Pilih
$a$ dengan $\varphi(a) = 1$; karena $a \notin \ker\varphi$ dan
kernelnya tertutup, ada bola $B(a, r)$ yang melewatkannya. Kini misalkan
$h \in E$ dengan $\varphi(h) \neq 0$: vektor $a - \frac{h}{\varphi(h)}$
berada di $\ker\varphi$, jadi di luar $B(a, r)$:
\[
\Bigl\Vert \frac{h}{\varphi(h)} \Bigr\Vert \geq r
\quad\Longrightarrow\quad
\abs{\varphi(h)} \leq \frac{\norm h}{r},
\]
sebuah ketaksamaan yang juga berlaku secara sepele ketika $\varphi(h) =
0$: yakni batas (4) pada \cref{thm:b2:nvs:continuouslinear}, jadi \omterm{def:b2:metric:continuity}{kontinu}.
\end{solution}

\begin{solution}{exo:b2:nvs:7}
Misalkan $f_n$ bernilai $0$ pada $\intcc{0}{\frac12 - \frac1n}$, afin
sampai bernilai $1$ di $\frac12$, lalu $1$ pada $\intcc{\frac12}{1}$.
Untuk $m \geq n$, $f_m - f_n$ bertumpu pada interval berpanjang
$\frac1n$ dengan nilai di $\intcc{-1}{1}$: jadi $\norm{f_m - f_n}_1 \leq
\frac1n$, sehingga Cauchy.

Andaikan $f_n \to f$ terhadap $\norm\cdot_1$ dengan $f$ \omterm{def:b2:metric:continuity}{kontinu}. Pada
$\intcc{0}{\frac12 - \delta}$ (dengan $\delta$ tetap): $\int \abs{f} =
\int\abs{f - f_n} \leq \norm{f - f_n}_1 \to 0$ untuk $n >
\frac1\delta$, jadi $\int_0^{1/2 - \delta}\abs f = 0$, dan berkat
kepositifan tegasnya $f = 0$ di sana --- untuk setiap $\delta$: jadi $f =
0$ pada $\intoo{0}{\frac12}$. Serupa itu $f = 1$ pada
$\intcc{\frac12}{1}$ (sebab semua $f_n$ bernilai $1$ di sana). Berkat
kekontinuan di $\frac12$: $0 = 1$, mustahil. Jadi tak ada limitnya:
ruangnya tidak \omterm{def:b2:metric:complete}{lengkap}.
\end{solution}

\begin{solution}{exo:b2:nvs:8}
Terdefinisi dengan baik dan \omterm{def:b2:metric:continuity}{kontinu}: $\abs{\varphi(f)} \leq \sum 2^{-n}
\norm f_\infty = \norm f_\infty$, jadi $\varphi$ merupakan bentuk linear
dengan $\vertiii\varphi \leq 1$ (deretnya konvergen mutlak untuk
tiap $f$).

\omterm{def:b2:nvs:norm}{Normanya} $1$: tetapkan $K$; titik $1, \frac12, \dots, \frac1K$ berbeda
sepasang demi sepasang, jadi ada $f_K$ yang \omterm{def:b2:metric:continuity}{kontinu} dengan
$\norm{f_K}_\infty \leq 1$ dan $f_K(\frac1n) = (-1)^n$ untuk $n \leq K$
(lewat interpolasi afin sepotong-sepotong, dan konstan di dekat $0$). Maka
\[
\varphi(f_K) \geq \sum_{n=1}^{K} 2^{-n} - \sum_{n > K} 2^{-n}
= 1 - 2^{-K+1} \xrightarrow[K \to \infty]{} 1 .
\]

Tidak tercapai: jika $\norm f_\infty \leq 1$ dan $\varphi(f) = 1$, maka
tiap sukunya wajib menyumbang maksimumnya: $(-1)^n f(\frac1n) = 1$ untuk
setiap $n$ (sebab jika tidak, kekurangan tegas satu suku tak dapat
ditutup, karena semua sukunya $\leq 2^{-n}$). Jadi $f(\frac1n) =
(-1)^n$; padahal $\frac1n \to 0$ dan $f$ \omterm{def:b2:metric:continuity}{kontinu} di $0$, sehingga
kekonvergenan $(-1)^n$ yang bertentangan itu terpaksa terjadi. Karenanya
supremumnya bukan maksimum --- dan hal itu mustahil dalam dimensi
hingga, tempat bola satuan tertutupnya \omterm{def:b2:metric:compact}{kompak}.
\end{solution}

\begin{solution}{exo:b2:nvs:9}
\emph{Barisan terbatas:} barisan terbatas berada di suatu bola
tertutup $\overline B(0, R) = R\,\overline B(0,1)$, yang \omterm{def:b2:metric:compact}{kompak} (sebab ia
peta bola satuan yang \omterm{def:b2:metric:compact}{kompak} di bawah homeomorfisma $x \mapsto Rx$):
jadi sarikan saja di sana.

\emph{Bentuk linear:} jika $\dim E < \infty$, setiap pemetaan linear dari
$E$ bersifat \omterm{def:b2:metric:continuity}{kontinu} (\cref{thm:b2:nvs:finitedim}). Sebaliknya,
andaikan $\dim E = \infty$ (yang sebenarnya, menurut
\cref{thm:b2:nvs:riesz}, sudah disingkirkan oleh hipotesis
kekompakannya --- inti pertanyaan ini adalah implikasi antara kedua
sifat itu pada ruang bernorma umum): pilihlah barisan $(e_n)$ yang bebas
linear dan ternormalkan, lengkapi menjadi basis aljabar (yang diterima
saja), lalu tetapkan $\varphi(e_n) = n$ dan $\varphi = 0$ pada vektor
basis lainnya, diperluas secara linear. Maka $\abs{\varphi(e_n)} = n$
dengan $\norm{e_n} = 1$: jadi ia tak terbatas pada bola satuannya, yakni
tidak \omterm{def:b2:metric:continuity}{kontinu}. Karenanya ``semua bentuk \omterm{def:b2:metric:continuity}{kontinu}'' memaksa dimensi hingga.
\end{solution}

\begin{solution}{exo:b2:nvs:10}
Cauchy--Schwarz terhadap fungsi konstan $1$: $\norm f_1 =
\int_0^1 \abs f\cdot 1 \leq \bigl(\int_0^1
f^2\bigr)^{1/2}\bigl(\int_0^1 1\bigr)^{1/2} = \norm f_2$. Dan
$\norm f_2^2 = \int f^2 \leq \norm f_\infty^2$. Untuk $f_n(x) = x^n$:
\[
\norm{f_n}_1 = \frac1{n+1}, \qquad
\norm{f_n}_2 = \frac1{\sqrt{2n+1}}, \qquad
\norm{f_n}_\infty = 1 .
\]
Maka $\norm{f_n}_2/\norm{f_n}_1 = \frac{n+1}{\sqrt{2n+1}} \to
\infty$ dan $\norm{f_n}_\infty/\norm{f_n}_2 = \sqrt{2n+1} \to
\infty$: jadi tak ada ketaksamaan sebaliknya, dan tak ada pasangan yang \omterm{def:b2:nvs:equivalent}{setara}.
\end{solution}

\begin{solution}{exo:b2:nvs:11}
\emph{Batas bawah bagi jaraknya:} untuk $h \in \ker\varphi$ berlaku
$\abs{\varphi(x)} = \abs{\varphi(x - h)} \leq
\vertiii\varphi\,\norm{x - h}$; ambil infimumnya atas $h$:
$d(x, \ker\varphi) \geq \abs{\varphi(x)}/\vertiii\varphi$.

\emph{Batas atasnya:} kita boleh mengandaikan $\varphi(x) \neq 0$.
Diberikan $\varepsilon > 0$, pilih vektor satuan $u$ dengan
$\abs{\varphi(u)} \geq \vertiii\varphi - \varepsilon > 0$ lalu tetapkan
$h = x - \frac{\varphi(x)}{\varphi(u)}\,u$: maka $\varphi(h) = 0$ dan
\[
\norm{x - h} = \frac{\abs{\varphi(x)}}{\abs{\varphi(u)}}
\leq \frac{\abs{\varphi(x)}}{\vertiii\varphi - \varepsilon}.
\]
Lalu ambil $\varepsilon \to 0$: $d(x, \ker\varphi) \leq
\abs{\varphi(x)}/\vertiii\varphi$, jadi keduanya sama.

\emph{Ketaktercapaiannya:} ambil $\varphi$ dari \cref{exo:b2:nvs:8}
($\vertiii\varphi = 1$, yang tidak tercapai) dan sembarang $x$ dengan
$\varphi(x) \neq 0$. Seandainya suatu $h \in \ker\varphi$ mewujudkan
$\norm{x - h} = \abs{\varphi(x)}$, maka vektor satuan $v = (x -
h)/\norm{x - h}$ akan memenuhi $\abs{\varphi(v)} =
\abs{\varphi(x)}/\norm{x - h} = 1 = \vertiii\varphi$: jadi
\omterm{thm:b2:nvs:continuouslinear}{norma operatornya} tercapai --- kontradiksi.
\end{solution}

\begin{solution}{exo:b2:nvs:12}
Ketaksamaan $N_1 \leq N_2$ tak lain kemonotonan supremum terhadap
daerahnya, jadi identitas $(E, N_2) \to (E, N_1)$ bersifat
Lipschitz-$1$. Untuk $P_n(x) = (x/2)^n$: $N_2(P_n) = 1$ (tercapai di $x =
2$) sedangkan $N_1(P_n) = 2^{-n}$: jadi batas $N_2 \leq CN_1$ akan
memberi $1 \leq C2^{-n}$ untuk setiap $n$: mustahil. Karenanya kedua
\omterm{def:b2:nvs:norm}{normanya} tak \omterm{def:b2:nvs:equivalent}{setara} dan identitas sebaliknya merupakan bijeksi linear
yang tidak \omterm{def:b2:metric:continuity}{kontinu}.

\emph{Ketaklengkapannya:} misalkan $S_n = \sum_{k=0}^{n}\frac{X^k}{k!}$.
Untuk $m > n$, $N_1(S_m - S_n) \leq \sum_{k>n}\frac1{k!} \to 0$: jadi
Cauchy terhadap $N_1$. Jika $S_n \to P$ di $(E, N_1)$, maka titik demi
titik $P(x) = \lim S_n(x) = \eu^x$ pada $\intcc{0}{1}$; padahal
polinomial berderajat $d$ tak mungkin sama dengan $\eu^x$ pada sebuah
interval (turunkan $d + 1$ kali: ruas kirinya mati, $\eu^x$ tidak). Jadi
tak ada limitnya di $E$: tidak \omterm{def:b2:metric:complete}{lengkap}.

Dalam dimensi hingga ketiga gejala itu mustahil: semua \omterm{def:b2:nvs:equivalent}{norma
setara}, setiap ruang bernorma bersifat \omterm{def:b2:metric:complete}{lengkap}, dan invers
sebuah bijeksi linear bersifat linear dari ruang berdimensi hingga,
jadi \omterm{def:b2:metric:continuity}{kontinu} (\cref{thm:b2:nvs:finitedim}).
\end{solution}

\begin{solution}{pb:b2:nvs:1}
\textbf{1.} Calon yang patut ditinjau membentuk $K = \{f \in F
: \norm{x - f} \leq \norm x\}$: tak kosong (sebab $0 \in K$), tertutup
(sebab ia prapeta sebuah interval tertutup di bawah pemetaan \omterm{def:b2:metric:continuity}{kontinu} $f
\mapsto \norm{x - f}$, yang diiriskan dengan $F$ yang tertutup,
\cref{cor:b2:nvs:closedsubspace}), dan terbatas (sebab $\norm f \leq
\norm{f - x} + \norm x \leq 2\norm x$). Di dalam $F$ yang
berdimensi hingga, tertutup dan terbatas berarti \omterm{def:b2:metric:compact}{kompak}
(\cref{thm:b2:nvs:finitedim}); lalu fungsi \omterm{def:b2:metric:continuity}{kontinu} $f
\mapsto \norm{x - f}$ mencapai infimumnya pada $K$, dan infimum itu sama
dengan infimum atas seluruh $F$ (sebab sembarang $f \notin K$ memberi
$\norm{x - f} > \norm x \geq \inf$).

\textbf{2.} Kesamaan jajargenjang di $(\R^n,
\norm\cdot_2)$: $\norm{u + v}^2 + \norm{u - v}^2 = 2\norm u^2 +
2\norm v^2$ (uraikan kuadrat jumlah koordinatnya). Untuk
$u \neq v$ yang satuan:
\[
\Bigl\Vert\frac{u+v}2\Bigr\Vert^2 = 1 - \frac{\norm{u -
v}^2}{4} < 1 .
\]
Tidak cembung tegas: untuk $\norm\cdot_\infty$, ambil $u = (1, 1, 0,
\dots)$, $v = (1, -1, 0, \dots)$: keduanya vektor satuan dengan titik
tengah $(1, 0, \dots)$ yang bernorma $1$; sedangkan untuk
$\norm\cdot_1$, ambil $u = (1, 0, \dots)$, $v = (0, 1, 0, \dots)$: titik
tengahnya $(\frac12, \frac12, 0, \dots)$ yang bernorma $1$.

\textbf{3.} Misalkan $d = d(x, F)$. Jika $d = 0$: maka $x \in \overline F
= F$ dan satu-satunya pemberi minimumnya adalah $x$. Jika $d > 0$ dan
$f_1 \neq f_2$ keduanya memberi minimum: maka $u = \frac{x - f_1}{d}$ dan
$v = \frac{x - f_2}{d}$ merupakan vektor satuan yang berbeda, sehingga
\[
\Bigl\Vert x - \frac{f_1 + f_2}2\Bigr\Vert
= d\,\Bigl\Vert\frac{u + v}2\Bigr\Vert < d ,
\]
dengan $\frac{f_1 + f_2}2 \in F$: dan itu bertentangan dengan definisi
$d$. Jadi pemberi minimumnya tunggal.

\textbf{4.} Di sini $\norm{(0,1) - t(1,0)}_\infty = \max(\abs t, 1)
\geq 1$, dengan kesamaan bila dan hanya bila $\abs t \leq 1$: jadi
pemberi minimumnya membentuk ruas $\{t(1, 0) : t \in \intcc{-1}{1}\}$,
semuanya berjarak $1$ --- ketunggalannya gagal justru karena bola
perseginya bersisi datar (pertanyaan 2).

\textbf{5.} Misalkan $M = \max f$ dan $m = \min f$ (yang tercapai berkat
kekompakan). Untuk sembarang konstanta $c$: $\sup\abs{f - c} \geq
\max(M - c,\, c - m) \geq \frac{M - m}2$, dengan ketaksamaan terakhirnya
berlaku karena kedua besaran itu berata-rata $\frac{M-m}2$; dan kesamaan
pada keduanya memaksa $M - c = c - m$, yakni $c = c^* = \frac{M + m}2$.
Sebaliknya $\sup\abs{f - c^*} = \max(M - c^*, c^* - m) =
\frac{M - m}2$. Jadi konstanta terbaiknya tunggal. Untuk $f(x) = x^2$ pada
$\intcc01$: $c^* = \frac12$, dengan jarak $\frac12$.

\textbf{6.} Dari $\cos(n{+}1)\theta + \cos(n{-}1)\theta =
2\cos\theta\cos n\theta$: polinomial yang ditetapkan oleh $T_0 = 1$,
$T_1 = X$, $T_{n+1} = 2XT_n - T_{n-1}$ memenuhi $T_n(\cos\theta)
= \cos n\theta$ secara induktif. Ketunggalannya: dua polinomial yang
berimpit pada $\intcc{-1}{1}$ (tak hingga banyak titik) pastilah sama.
Sekali lagi secara induktif: $\deg T_n = n$ dengan koefisien utama
$2^{n-1}$ untuk $n \geq 1$ (sebab $T_1$ berkoefisien $1 = 2^0$, dan
rekurensinya melipatduakannya).

\textbf{7.} Setiap $x \in \intcc{-1}{1}$ berbentuk $\cos\theta$, dan
$\abs{\cos n\theta} \leq 1$. Di $\eta_k = \cos\frac{k\pi}n$:
$T_n(\eta_k) = \cos k\pi = (-1)^k$, dan $1 = \eta_0 > \eta_1 >
\dots > \eta_n = -1$. Akarnya: $\cos n\theta = 0$ bila dan hanya bila
$\theta = \frac{(2k-1)\pi}{2n}$, yakni $n$ titik berbeda
$\cos\frac{(2k-1)\pi}{2n}$; dan karena $\frac{(k-1)\pi}n <
\frac{(2k-1)\pi}{2n} < \frac{k\pi}n$, tiap akarnya terletak tegas
di antara dua ekstremum yang berurutan.

\textbf{8.} $T_2 = 2X^2 - 1$, $T_3 = 4X^3 - 3X$, $T_4 = 8X^4 -
8X^2 + 1$. Untuk $T_3$: $T_3(1) = 1$, $T_3(\tfrac12) = \tfrac12 -
\tfrac32 = -1$, $T_3(-\tfrac12) = 1$, $T_3(-1) = -1$: pergantian
tandanya sempurna.

\textbf{9.} Misalkan $u_\pm = x \pm \sqrt{x^2 - 1}$ untuk $x \geq 1$:
yakni akar $z^2 - 2xz + 1$, dengan $u_+u_- = 1$. Barisan
$s_n = \frac{u_+^n + u_-^n}2$ memenuhi $s_{n+1} = 2x\,s_n -
s_{n-1}$ (rekurensi bergaya Newton dari persamaan kuadratnya), $s_0 = 1$,
$s_1 = x$: yakni rekurensi dan nilai awal yang sama seperti $n \mapsto
T_n(x)$, jadi $s_n = T_n(x)$ untuk setiap $n$. Karena $0 < u_- \leq 1
\leq u_+$ dengan $u_+ > 1$ untuk $x > 1$: maka $T_n(x) \geq
\frac{u_+^n}2 \to \infty$ dan $T_n(x) \sim \frac12\bigl(x +
\sqrt{x^2-1}\bigr)^n$. (Untuk $x \leq -1$ pakailah paritas $T_n(-x)
= (-1)^nT_n(x)$, yang jelas dari rekurensinya.)

\textbf{10.} Untuk setiap $\theta$: $T_m\bigl(T_n(\cos\theta)\bigr)
= T_m(\cos n\theta) = \cos mn\theta =
T_{mn}(\cos\theta)$. Jadi polinomial $T_m \circ T_n$ dan
$T_{mn}$ berimpit pada $\intcc{-1}{1}$, sehingga keduanya sama.

\textbf{11.} Polinomial $D = Q_n - P$ berderajat $\leq n - 1$ (sebab suku
utamanya yang monik saling coret). Di ekstremumnya: $(-1)^kD(\eta_k) =
2^{1-n} - (-1)^kP(\eta_k) \geq 2^{1-n} - \abs{P(\eta_k)} > 0$ menurut
hipotesisnya. Jadi $D$ mengambil nilai tak nol bertanda berselang-seling
di $n + 1$ titik menurun $\eta_0 > \dots > \eta_n$: sehingga menurut
teorema nilai antara ia punya sedikitnya $n$ akar berbeda,
satu di tiap interval \omterm{def:b2:metric:topology}{terbuka} $\intoo{\eta_{k}}{\eta_{k-1}}$. Padahal
polinomial tak nol yang berderajat $\leq n - 1$ tak mungkin punya $n$ akar;
sedangkan $D = 0$ bertentangan dengan tanda tegas tadi. Jadi
kontradiksi: $\sup\abs P \geq 2^{1-n}$ untuk setiap $P$ monik berderajat $n$.

\textbf{12.} Kini $(-1)^kD(\eta_k) = 2^{1-n} - (-1)^kP(\eta_k)
\geq 0$ sebab $\abs{P} \leq 2^{1-n}$. Pada tiap
$\intcc{\eta_k}{\eta_{k-1}}$ ($k = 1, \dots, n$), nilai $D$ di kedua
ujungnya bertanda lemah yang berlawanan: jadi teorema nilai antara
menghasilkan sebuah nol $z_k$ di interval tertutup itu. Jika $z_k$ itu
dapat dipilih berbeda sepasang demi sepasang, maka $D \neq 0$ yang
berderajat $\leq n-1$ punya $n$ akar: kontradiksi. Dua interval berurutan
hanya dapat berbagi nol $z_k = z_{k+1} = \eta_k$ dengan $0 < k < n$
(yakni titik dalam). Di sana, $D(\eta_k) = 0$ berarti $P(\eta_k) =
(-1)^k2^{1-n}$, yaitu nilai ekstremal $P$ pada $\intcc{-1}{1}$ yang
tercapai di titik \emph{dalam}: jadi $P'(\eta_k) = 0$; dan
$\eta_k$ juga ekstremum dalam bagi $T_n$: jadi $Q_n'(\eta_k) =
0$. Karenanya $D'(\eta_k) = 0$: sehingga $\eta_k$ akar bermultiplisitas
$\geq 2$, yang menutupi interval yang dibagi tadi. Pada semua kasusnya
$D$ punya sedikitnya $n$ akar dihitung dengan multiplisitasnya, padahal
derajatnya $\leq n - 1$, jadi $D = 0$: yakni $P = Q_n$. Teorema ekstremal
Chebyshev pun terbukti: pemberi minimum monik yang tunggal adalah
$2^{1-n}T_n$, bernorma supremum $2^{1-n}$.

\textbf{13.} Polinomial monik berderajat $n$ persis semua
$X^n - R$ dengan $R \in \R_{n-1}[X]$, jadi
\[
d_\infty\bigl(X^n, \R_{n-1}[X]\bigr) = \min_{P \text{
monik}}\ \sup_{\intcc{-1}{1}}\abs P = 2^{1-n},
\]
yang tunggal di $R^* = X^n - Q_n$. Dengan penyulihan $x = \frac{1+t}2$:
jika $P$ monik berderajat $n$ pada $\intcc01$, maka $t \mapsto
2^nP\bigl(\frac{1+t}2\bigr)$ monik pada $\intcc{-1}{1}$ dengan
supremum sama dengan $2^n\sup_{\intcc01}\abs P$: karenanya
$\sup_{\intcc01}\abs P \geq 2^{-n}\cdot2^{1-n} = 2^{1-2n}$, dengan
kesamaan persis untuk $P^*(x) = 2^{-n}Q_n(2x - 1)$: jadi pada
$\intcc{0}{1}$ jaraknya $2^{1-2n}$.

\textbf{14.} Polinomial $\omega$ monik berderajat $n$, jadi
$\sup_{\intcc{-1}{1}}\abs\omega \geq 2^{1-n}$ dengan kesamaan bila dan
hanya bila $\omega = Q_n = 2^{1-n}T_n$, yakni bila dan hanya bila
simpulnya adalah $n$ akar $T_n$: $x_k = \cos\frac{(2k-1)\pi}{2n}$ ---
yaitu \emph{simpul Chebyshev}. Nilai minimumnya: $2^{1-n}$. Simpul yang
berjarak sama justru lebih buruk; faktor galat interpolasinya diminimumkan
dengan merapatkan simpul di dekat ujungnya.

\textbf{15.} Kasus $n = 2$ dengan tangan: $x^2$ menjelajahi $\intcc01$,
jadi $\sup_{\intcc{-1}{1}}\abs{x^2 - c} = \max(\abs c, \abs{1 - c})
\geq \frac12$, yang minimum di $c = \frac12$: jadi polinomial kuadrat
monik yang minimal adalah $x^2 - \frac12 = \frac12(2x^2 - 1) = Q_2$,
bernilai $\frac12 = 2^{1-2}$. Untuk $n = 10$ pada $\intcc{0}{1}$: $2^{1-20} =
2^{-19} \approx 1.9\cdot10^{-6}$. Jadi ada polinomial berderajat $9$
yang tetap berjarak dua per sejuta dari $x^{10}$ pada seluruh
$\intcc01$: pada skala itu kedua grafiknya tak terbedakan
--- kedataran $x^{10}$ di dekat $0$ membiarkan derajat yang lebih rendah
mengerjakan segalanya.

\textbf{16.} Di sini $\norm{g_k}_\infty = g_k(1) = 1$. Untuk $j > k$
tetapkan $a = 2^k$, $b = 2^j \geq 2a$, lalu nilailah di $x_0 = 2^{-1/a}$
(sehingga $x_0^a = \frac12$):
\[
g_k(x_0) - g_j(x_0) = \frac12 - \Bigl(\frac12\Bigr)^{b/a}
\geq \frac12 - \frac14 = \frac14 .
\]
Jadi $\norm{g_k - g_j}_\infty \geq \frac14$ untuk setiap $j \neq k$:
sehingga tak ada barisan bagian yang Cauchy, dan tak satu pun konvergen.
Bola satuan tertutup $\bigl(C(\intcc01), \norm\cdot_\infty\bigr)$ tidak \omterm{def:b2:metric:compact}{kompak}.

\textbf{17.} Subruang $F$ tertutup (\cref{cor:b2:nvs:closedsubspace})
dan sejati: pilih $y \notin F$, sehingga $\delta = d(y, F) > 0$. Menurut
pertanyaan 1 jaraknya tercapai di suatu $f^* \in F$. Tetapkan $x
= \frac{y - f^*}{\delta}$, sebuah vektor satuan (sebab $\norm{y - f^*} =
\delta$). Untuk setiap $g \in F$:
\[
\norm{x - g} = \frac{\norm{y - (f^* + \delta g)}}{\delta}
\geq \frac{\delta}{\delta} = 1 ,
\]
sebab $f^* + \delta g \in F$. Karenanya $d(x, F) \geq 1$; sedangkan $d(x,
F) \leq \norm{x - 0} = 1$: jadi tepat $1$.

\textbf{18.} Pada $E$ yang berdimensi tak hingga, bangunlah vektor satuan
secara induktif: $x_1$ sembarang; lalu diberikan $x_1, \dots, x_k$,
subruang $F_k = \operatorname{Vect}(x_1, \dots, x_k)$ berdimensi
hingga, jadi sejati, sehingga pertanyaan 17 memberi vektor satuan
$x_{k+1}$ dengan $d(x_{k+1}, F_k) = 1$: khususnya
$\norm{x_{k+1} - x_i} \geq 1$ untuk $i \leq k$. Barisan itu berjarak
antara $\geq 1$: jadi bola satuannya memuat barisan
tanpa barisan bagian yang konvergen, sehingga ia tidak \omterm{def:b2:metric:compact}{kompak} ---
itulah teorema Riesz, dengan konstanta tajam $1$.

\textbf{19.} Misalkan $h_n$ afin pada tiap separuh
$\bigl[\frac1{n+1}, \frac1n\bigr]$, naik dari $0$ ke $1$ di
titik tengahnya lalu turun kembali ke $0$, dan nol di tempat lain: ia
\omterm{def:b2:metric:continuity}{kontinu} dengan $\norm{h_n}_\infty = 1$. Untuk $n \neq m$ penyangganya
bertemu paling banyak di satu ujung bersama, tempat keduanya nol;
sedangkan di puncak $h_n$ berlaku $h_m = 0$: jadi $\norm{h_n -
h_m}_\infty = 1$ persis. Untuk $x > 0$ yang tetap: $h_n(x) = 0$ begitu
$\frac1n < x$, dan $h_n(0) = 0$ selalu: jadi $h_n \to 0$ titik demi
titik; tetapi $\norm{h_n - 0}_\infty = 1$: jadi tidak seragam. Itulah
gugusan gamblang berjarak antara $1$ di dalam bola satuannya.

\textbf{20.} Bola berjari-jari $\frac13$ berdiameter $\leq
\frac23 < 1$, jadi ia memuat paling banyak satu $h_n$ (sebab dua di
antaranya berjarak $1$). Berhingga banyak bola semacam itu memuat
berhingga banyak dari tak hingga banyak $h_n$: jadi keduanya tak dapat
menyelimuti bola satuannya. Keterbatasan total --- yang dinikmati \omterm{def:b2:metric:compact}{ruang
metrik kompak}, menurut bukti \cref{thm:b2:metric:borellebesgue} ---
gagal separah-parahnya.

\textbf{21.} Jika $Ax = \lambda x$ dengan $x \neq 0$, maka
$\abs\lambda\,\norm x = \norm{Ax} \leq \vertiii A\,\norm x$,
sehingga $\abs\lambda \leq \vertiii A$. Untuk
$\begin{psmallmatrix}1 & -2\\ 3 & 1\end{psmallmatrix}$:
$\vertiii A_\infty = \max(1 + 2,\ 3 + 1) = 4$
(\cref{exo:b2:nvs:4}), jadi setiap \omterm{def:b2:reduction:eigen}{nilai eigennya} bermodulus $\leq
4$; padahal $\chi_A = X^2 - 2X + 7$ memberi $\lambda = 1 \pm
\iu\sqrt6$ yang bermodulus $\sqrt7 \approx 2.65$: jadi batasnya
sahih tetapi tidak tajam.

\textbf{22.} $N_P$ memang \omterm{def:b2:nvs:norm}{norma}: $N_P(x) = 0$ memaksa $P^{-1}x = 0$,
jadi $x = 0$; sedangkan kehomogenan dan ketaksamaan segitiganya diwarisi
dari $\norm\cdot_\infty$ lewat $P^{-1}$ yang linear. \omterm{thm:b2:nvs:continuouslinear}{Norma
operatornya}: dengan $y = P^{-1}x$ dan $D = \mathrm{diag}(\lambda_i)$,
\[
N_P(Ax) = \norm{P^{-1}AP\,y}_\infty = \norm{Dy}_\infty,
\]
jadi $\vertiii A_{N_P}$ tak lain \omterm{thm:b2:nvs:continuouslinear}{norma operator} terhadap
$\norm\cdot_\infty$ bagi $D$, yakni jumlah mutlak barisnya yang terbesar
(\cref{exo:b2:nvs:4}): $\max_i\abs{\lambda_i}$.

\textbf{23.} Jika semua $\abs{\lambda_i} < 1$: maka $N_P(A^kx) \leq
\rho^kN_P(x)$ dengan $\rho = \max\abs{\lambda_i} < 1$, jadi $A^kx
\to 0$ untuk setiap $x$, sehingga $A^k \to 0$ pada \omterm{def:b2:nvs:norm}{norma} apa pun di
$\mathcal{M}_n$ (semuanya \omterm{def:b2:nvs:equivalent}{setara} dalam dimensi hingga,
\cref{thm:b2:nvs:finitedim}; sebab kekonvergenan $A^ke_j$ untuk tiap
$j$ tak lain kekonvergenan unsur demi unsur). Jika ada $\abs{\lambda}
\geq 1$ dengan \omterm{def:b2:reduction:eigen}{vektor eigen} $x$: maka $\norm{A^kx} =
\abs\lambda^k\norm x \not\to 0$. Untuk $A =
\frac14\begin{psmallmatrix}1 & 2\\ 2 & 1\end{psmallmatrix}$: \omterm{def:b2:reduction:eigen}{nilai
eigennya} $\frac14(1 \pm 2) = \frac34, -\frac14$, keduanya bermodulus
$< 1$: jadi $A^k \to 0$.

\textbf{24.} Keduanya \omterm{def:b2:nvs:norm}{norma} pada $\R_n[X]$ yang berdimensi
hingga, jadi \omterm{def:b2:nvs:equivalent}{setara} untuk tiap $n$. Ambil polinomial monik
minimal $P^*_n$ pada pertanyaan 13 di $\intcc01$: koefisien
$X^n$-nya $1$, jadi $N_c(P^*_n) \geq 1$, sedangkan
$\norm{P^*_n}_{\intcc01} = 2^{1-2n}$. Karenanya
\[
C_n \geq \frac{N_c(P^*_n)}{\norm{P^*_n}_{\intcc01}}
\geq 2^{2n-1} .
\]
Konstanta kesetaraannya meledak seiring dimensinya: pada
gabungan $\R[X]$ tak ada satu pun konstanta yang melayani, dan itulah
persis ketaksetaraan yang terlihat pada \cref{exo:b2:nvs:12} --- ``semua
\omterm{def:b2:nvs:equivalent}{norma setara}'' adalah teorema tentang satu dimensi setiap kali,
dan dimensi tak hingga adalah tempat ia mati.

\textbf{25.} Kekompakan himpunan tertutup dan terbatas dalam dimensi
hingga menghasilkan keberadaan hampiran terbaik (pertanyaan
1) serta menggerakkan pencacahan nol di $n + 1$ titik ekstremalnya
(pertanyaan 11--12, lewat supremum yang tercapai). Kecembungan tegas
menguasai ketunggalan hampiran terbaik --- bola yang bulat memberi satu
pemberi minimum, dan bola bersisi datar memberi seruas penuh (pertanyaan
2--4). Gugusan vektor satuan berjarak antara
$1$ (pertanyaan 17--19) menghancurkan kekompakan bola satuannya
sekaligus keterbatasan totalnya (pertanyaan 20). Pertanyaan 24
mengukur keruntuhannya: konstanta yang menghubungkan dua \omterm{def:b2:nvs:norm}{norma} pada
$\R_n[X]$ tumbuh seperti $4^n$, jadi tak ada perbandingan seragam yang
selamat sampai ke $\R[X]$. Puncaknya adalah teorema ekstremal
Chebyshev (pertanyaan 11--12): pemberi minimum monik yang tunggal
$2^{1-n}T_n$. Adapun hampiran terbaik menemukan rumah modernnya di
ruang Hilbert, tempat teorema proyeksi menggantikan
kekompakan dengan kelengkapan ditambah kesamaan jajargenjang ---
yang dibuktikan secara jujur pada jilid Tahun ke-3.

\end{solution}
